\chapter*{常见缩写和术语先导}
\addcontentsline{toc}{chapter}{常见缩写和术语先导}

第一册会把 C++ 程序一路追到二进制、汇编、进程和测量证据。下面这些缩写和基础术语会反复出现，先读一遍能减少正文中的断点。

\begin{description}
  \item[CPU] Central Processing Unit，中央处理器。本书关注它如何取指、译码、执行指令，以及这些事实如何影响 C++ 程序。
  \item[ISA] Instruction Set Architecture，指令集架构。它规定软件能看见哪些指令、寄存器和状态；x86-64 就是一种 ISA。
  \item[x86-64/AMD64/IA-32/RISC-V] x86-64 和 AMD64 常指同一类 64 位 x86 架构；IA-32 是 32 位 x86 架构名；RISC-V 是另一类开放 ISA。ABI、寄存器和调用约定不能在这些目标之间直接套用。
  \item[OS/PID/PC/VM] OS 是 Operating System，操作系统；PID 是 process id，进程编号；PC 是 program counter，程序计数器；VM 在本书中多指 virtual memory，虚拟内存。第一册看程序运行时，会频繁通过 PID 进入 \code{/proc}、调试器和系统调用证据。
  \item[PR/ID/UNIX] PR 是 pull request，代码合并请求；ID 是 identifier 或 identity，表示身份编号；UNIX 是一类操作系统传统和接口生态，Linux 继承了大量 UNIX/POSIX 思想。
  \item[PROC] Linux 中常指 \code{/proc} 伪文件系统。它把进程、内核和硬件状态暴露成可读文件，是本书做运行时证据采集的入口之一。
  \item[EOF/RLIMIT/RSS] EOF 是 End Of File，输入结束；RLIMIT 是进程资源限制；RSS 是 Resident Set Size，常驻内存大小。它们经常出现在进程、文件和资源排查里。
  \item[CC/CXX/CFLAGS] 构建环境变量。CC 常指定 C 编译器，CXX 常指定 C++ 编译器，CFLAGS/CXXFLAGS 常指定编译选项；实验报告要记录它们是否影响结果。
  \item[C++20] C++ 标准版本名。第一册示例会使用 \code{std::bit_cast}、benchmark harness 等 C++20 能力，读代码时要先确认编译选项是否启用这个标准。
  \item[ABI] Application Binary Interface，应用二进制接口。它规定函数调用、寄存器使用、栈布局、参数传递和符号链接边界。
  \item[API] Application Programming Interface，源码层接口；ABI 是二进制层接口。源代码能编译不代表 ABI 一定兼容。
  \item[CI/I/O] CI 是 Continuous Integration，持续集成；I/O 是 Input/Output，输入输出。它们会进入实验可信度：自动检查能否重放，计时是否混入文件或终端输出。
  \item[STL] Standard Template Library，C++ 标准库容器、迭代器和算法的常见统称。ABI 章节会提醒：不要把 STL 容器当长期稳定的二进制接口直接暴露。
  \item[O0/O2/O3/DNDEBUG] 常见构建选项。O0 几乎不优化，方便调试；O2/O3 会启用更多优化；DNDEBUG 常用于关闭 \code{assert}。
  \item[IR] Intermediate Representation，中间表示。编译器把源码翻译成 IR 后做优化，再 lowering 到机器相关表示。
  \item[AST] Abstract Syntax Tree，抽象语法树。它表示源码语法结构，例如一个函数、表达式或声明之间的层级关系。
  \item[CFG] Control-Flow Graph，控制流图。它把基本块和跳转边连起来，帮助解释分支、循环和优化。
  \item[LLVM/GCC/Clang] LLVM 是编译器基础设施项目；GCC 是 GNU Compiler Collection；Clang 是 LLVM 生态中的 C/C++ 前端。本书用它们生成 IR、汇编和诊断证据。
  \item[SSA/GIMPLE/LTO] SSA 是 Static Single Assignment，静态单赋值形式；GIMPLE 是 GCC 的中间表示之一；LTO 是 Link-Time Optimization，链接期优化。
  \item[DCE/LICM/mem2reg] 常见优化名。DCE 删除无用计算；LICM 把循环不变计算外提；mem2reg 把可证明安全的栈槽提升为 SSA 寄存器值。
  \item[objdump/readelf/strace/perf] Linux/C++ 底层实验常用工具。objdump 看反汇编和 section，readelf 看 ELF 头、符号、重定位和动态段，strace 看系统调用，perf 看采样和 PMU 计数器。它们回答的是不同层的问题，不能互相替代。
  \item[ASM/MSVC/COFF/SEH] ASM 是 assembly，汇编；MSVC 是 Microsoft Visual C++ 工具链；COFF 是一类目标文件格式族；SEH 是 Windows 结构化异常处理。
  \item[GAS/MASM/FFI] GAS 是 GNU assembler，GNU 汇编器语法生态；MASM 是 Microsoft Macro Assembler；FFI 是 Foreign Function Interface，跨语言调用边界。
  \item[ELF] Executable and Linkable Format，Linux 常见的目标文件、共享库和可执行文件格式。
  \item[ELF64/INTERP/NEEDED/JUMP\_SLOT] ELF64 表示 64 位 ELF 文件；INTERP 记录动态加载器路径；NEEDED 记录运行时需要加载的共享库；JUMP\_SLOT 是动态函数调用绑定相关重定位类型。
  \item[PT\_LOAD/BIND\_NOW/LD\_BIND\_NOW] PT\_LOAD 是 ELF 程序头中需要装载进内存的段；BIND\_NOW/LD\_BIND\_NOW 表示尽早完成动态符号绑定，而不是第一次调用时再懒绑定。
  \item[PLT/GOT] PLT 是 Procedure Linkage Table，GOT 是 Global Offset Table。它们共同支持动态链接下的函数和全局对象寻址。
  \item[GOTPCREL] x86-64 ELF 重定位/寻址相关记号，常出现在位置无关代码访问 GOT 项时。
  \item[DWARF/GDB] DWARF 是常见调试信息格式；GDB 是 GNU 调试器。前者把源码位置、变量和类型写进二进制旁的信息里，后者读取这些信息来断点和查看栈帧。
  \item[ASan/UBSan/TSan] AddressSanitizer、UndefinedBehaviorSanitizer、ThreadSanitizer。它们把内存错误、未定义行为和数据竞争变成诊断证据，但会改变代码和内存布局。
  \item[CFI] Call Frame Information，调用帧信息。它帮助调试器、异常处理和采样器从当前栈帧回到调用者。
  \item[RPATH/RUNPATH] ELF 中影响动态库搜索路径的字段。看到它们时，要追问程序运行时到底从哪个目录加载共享库。
  \item[LD/LD\_LIBRARY\_PATH/LD\_DEBUG] 动态链接器和运行时库搜索路径相关概念。它们能改变程序实际加载的共享库；\code{LD_DEBUG} 可以让动态链接器输出库查找和符号绑定过程。
  \item[mmap/FD/CLOEXEC] mmap 把文件或匿名页映射进进程虚拟地址空间；FD 是 file descriptor，文件描述符；CLOEXEC 表示执行新程序时自动关闭该描述符，避免资源泄漏到子进程。
  \item[RIP/RFLAGS/EFLAGS/ZF/CF/SF/OF/PF] x86-64 架构状态里的寄存器和标志。RIP 保存下一条指令地址；RFLAGS/EFLAGS 保存条件标志；ZF 表示结果为零；CF 表示无符号进位或借位；SF 表示符号位；OF 表示有符号溢出；PF 表示低字节奇偶。
  \item[SIB/LEA] SIB 是 x86 复杂寻址编码字段，用 base、index 和 scale 组合出有效地址；LEA 是 load effective address，常把地址计算当成整数算术使用。
  \item[REX/VEX/EVEX] x86-64 和向量指令的扩展前缀。它们会改变寄存器编号、操作数宽度或向量能力。
  \item[XMM/YMM] x86 SIMD 向量寄存器名。XMM 是 128 位，YMM 是 256 位。
  \item[CPL/DPL/RPL] x86 权限检查里的特权级名词。CPL 是当前特权级，DPL 是描述符特权级，RPL 是 selector 中携带的请求特权级。
  \item[IDT/PTE] IDT 是中断描述符表，PTE 是页表项。前者帮助处理器找到异常入口，后者记录虚拟页映射和权限。
  \item[TSS/RSP0/IST] TSS 是任务状态段；RSP0 是从用户态切到内核态时常用的 ring 0 栈指针字段；IST 是中断栈表。它们会在异常和中断入口里出现。
  \item[MSR/IA32/IPI] MSR 是模型专用寄存器；IA32 是 x86 体系中沿用的寄存器命名前缀；IPI 是处理器间中断。系统调用入口和跨核事件会碰到它们。
  \item[\#PF/CET/NX/XD] \#PF 是 page fault，页故障异常；CET 是 Control-flow Enforcement Technology，控制流保护技术；NX/XD 表示不可执行页权限。它们分别连接内存权限和控制流安全。
  \item[SIGSEGV/SIGPIPE] 常见信号。SIGSEGV 通常表示非法内存访问，SIGPIPE 通常表示向已经关闭的管道或 socket 写入。
  \item[SIGTERM/SIGINT] 常见终止/中断信号。SIGTERM 是请求进程终止，SIGINT 通常来自终端 Ctrl-C。
  \item[SIGABRT/SIGALRM] 常见信号。SIGABRT 常由 \code{abort} 或断言失败触发；SIGALRM 来自定时器告警。
  \item[SIGKILL/SIGBUS/SIGILL] SIGKILL 是不能被捕获的强制终止信号；SIGBUS 常和 mmap 生命周期、设备或对齐问题有关；SIGILL 表示执行了非法或当前 CPU 不支持的指令。
  \item[SIGSTKSZ] POSIX 信号备用栈相关常量，用来给 \code{sigaltstack} 选择基础栈大小。它不是错误码，而是 signal 处理的工程边界。
  \item[SA\_RESTART/SA\_SIGINFO/SA\_ONSTACK] POSIX signal flag。SA\_RESTART 允许部分系统调用被信号打断后自动重启；SA\_SIGINFO 让 handler 接收更详细的信号信息；SA\_ONSTACK 让 handler 使用备用栈。
  \item[EINTR/ENOENT/EAGAIN/EMFILE] 常见 errno。EINTR 表示系统调用被信号中断；ENOENT 表示路径不存在；EAGAIN 表示资源暂时不可用；EMFILE 表示进程文件描述符用尽。
  \item[EFAULT/EPIPE/ENOMEM/EINVAL/EWOULDBLOCK] 常见 errno。EFAULT 通常表示用户态地址无效；EPIPE 表示管道或 socket 对端已关闭；ENOMEM 表示内存不足；EINVAL 表示参数无效；EWOULDBLOCK 常和非阻塞 I/O 暂时不能完成有关。
  \item[POSIX/VFS] POSIX 是 Unix-like 系统接口规范；VFS 是 Linux 虚拟文件系统层。第一册讨论系统调用时，会把接口规则和内核抽象分开。
  \item[futex/epoll] futex 是 Linux 用户态同步原语和内核等待队列之间的接口，高层 mutex、condition variable 在竞争时常会落到它；epoll 是 Linux I/O 事件通知机制，用来等待多个文件描述符的可读、可写或错误事件。它们都是系统边界，不是普通容器。
  \item[TLS] Thread-Local Storage，线程局部存储。每个线程拥有一份变量实例，访问方式会影响 ABI 和生成代码。
  \item[ASLR] Address Space Layout Randomization，地址空间布局随机化。系统每次运行时随机化关键地址，降低攻击者预测地址的能力。
  \item[PIE/PIC] PIE 是位置无关可执行文件，PIC 是位置无关代码。它们让代码可以装载到不同地址运行。
  \item[UAF/ROP/RTTI/LSDA] UAF 是 use-after-free，释放后继续使用；ROP 是 return-oriented programming，利用返回地址拼接已有指令片段；RTTI 是 C++ 运行期类型信息；LSDA 是异常展开用的语言特定数据区域。
  \item[RTL] Register Transfer Language，GCC 后端常见的中间表示层之一。它比 GIMPLE 更接近机器寄存器和指令选择。
  \item[ALU/ROB] ALU 是算术逻辑单元；ROB 是重排序缓冲区。第一册只建立入口，第二册会系统展开乱序执行资源。
  \item[BTB/RSB/RISC] BTB 是分支目标缓存；RSB 是返回栈缓冲；RISC 是精简指令集计算机风格的架构类别。它们属于机器模型入口，第二册会展开性能影响。
  \item[ILP] Instruction-Level Parallelism，指令级并行。多条互不依赖的指令可以重叠执行，后续性能章节会用它解释展开、多累加器和向量化。
  \item[SIMD/SSE/AVX/AVX2/FMA] SIMD 是一条指令处理多个数据 lane；SSE、AVX、AVX2 是 x86 向量扩展；FMA 是 fused multiply-add，一条指令完成乘加并减少中间舍入。
  \item[AVX-512/ZMM] AVX-512 是更宽的 x86 向量扩展；ZMM 是 512 位向量寄存器名。它更强，但也可能带来频率、代码体积和平台覆盖问题。
  \item[TLB] Translation Lookaside Buffer，地址翻译缓存。它缓存虚拟页到物理页的映射，TLB miss 会增加访存成本。
  \item[L1/L1D/L2/L3] CPU 缓存层级。L1 最小最快，L1D 特指一级数据缓存，L2 居中，L3 通常更大并可能被多个核心共享。
  \item[PMU/TSC] PMU 是性能监控单元，提供硬件计数器；TSC 是 Time Stamp Counter，常用于读取处理器时间戳。二者都能辅助测量，但含义不同。
  \item[IPC/CPI/CSV/GB/GFLOP/QPS] IPC 是每周期执行指令数；CPI 是每条指令平均周期数；CSV 是逗号分隔表格文本；GB 是十亿字节量级；GFLOP/s 是每秒十亿次浮点运算；QPS 是每秒请求数。指标要和单位、公式、原始数据一起看。
  \item[p50/p95/p99] 分位数指标。p50 是中位数，p95/p99 描述尾部样本；性能章节用它们说明延迟抖动和极慢请求，而不是只看平均值。
  \item[LLC] Last-Level Cache，最后一级缓存，通常指多个核心共享的较大缓存。
  \item[SMT/COW/RMW/JIT] SMT 是同时多线程；COW 是写时复制；RMW 是读改写原子序列；JIT 是 Just-In-Time，运行时即时生成或编译代码。
  \item[NUMA] Non-Uniform Memory Access，非统一内存访问。不同 CPU 节点访问不同内存节点的成本不同。
  \item[HPC/GEMM/VNNI/FTZ/DAZ/MXCSR] HPC 是高性能计算；GEMM 是矩阵乘矩阵；VNNI 是低精度点积相关向量指令能力；FTZ/DAZ 是浮点非正规数处理模式；MXCSR 是 x86 SSE/AVX 浮点控制状态寄存器。
  \item[INTEGER/SSEUP/MEMORY/X87] System V AMD64 ABI 聚合分类中的类别名。INTEGER 通常走通用寄存器，SSE/SSEUP 关联向量/浮点寄存器，MEMORY 表示走内存，X87 关联旧 x87 浮点规则。
  \item[IEEE/NaN/Inf/ULP/NX/BSS/VMA/CFA] IEEE 754 是常见浮点标准；NaN/Inf 是特殊浮点值；ULP 是浮点末位单位；NX 是不可执行页权限；BSS 是未初始化全局数据段；VMA 是虚拟内存区域；CFA 是调用帧地址。
  \item[ODR/IDE/README] ODR 是 One Definition Rule，单一定义规则；IDE 是集成开发环境；README 是项目说明文件。它们属于工程可复现性的基础词。
  \item[JSON/LCG/BLAS] JSON 是结构化文本格式；LCG 是线性同余伪随机生成器；BLAS 是基础线性代数库接口族。它们常出现在实验配置、可复现输入和高性能库环境中。
  \item[KV/TTFT] KV 常指 Transformer 的 Key/Value cache；TTFT 是 Time To First Token。第一册只作为后续册的运行时例子出现，第三册会系统展开。
  \item[WSL2/UI] WSL2 是 Windows Subsystem for Linux 2，Windows 上的 Linux 环境；UI 是用户界面。性能实验要说明是否在虚拟化、容器或图形工具环境里运行。
  \item[UB] Undefined Behavior，未定义行为。C++ 程序触发 UB 后，标准不再规定结果，编译器优化也可能放大问题。
  \item[RAII] Resource Acquisition Is Initialization，把资源生命周期绑定到对象生命周期的 C++ 惯用法。
  \item[DMA/DRAM] DMA 是 Direct Memory Access，设备直接读写内存；DRAM 是主存使用的动态随机访问存储器。
  \item[SRAM/RCU] SRAM 是静态随机访问存储器，常用于 CPU cache；RCU 是 Read-Copy-Update，Linux 内核常见的读多写少同步思想。
  \item[register/stack/heap/frame] register 是寄存器，CPU 内部的高速小存储；stack 是调用栈，保存返回地址、局部对象和调用现场；heap 是动态分配区域；frame 是一次函数调用在栈和调试信息里的那一帧。
  \item[kernel/syscall/interrupt/exception] kernel 是操作系统内核；syscall 是用户态请求内核服务的入口；interrupt 是外部事件打断当前执行；exception 是处理器因当前指令或状态触发的异常。
  \item[symbol/section/relocation] symbol 是链接器能识别的名字；section 是目标文件里的分区；relocation 是“这里的地址或偏移等链接时再填”的修正记录。只看汇编而不看 relocation，很多地址解释会错。
  \item[layout/padding/alignment] layout 是对象或数据在内存中的排列；padding 是编译器插入的填充字节；alignment 是地址对齐要求。结构体大小、ABI 和 cache 行都依赖它们。
  \item[baseline/reference/oracle/checksum] baseline 是可重复对比的基线实现；reference 是定义正确结果的清晰实现；oracle 是比较结果是否满足合同的判定器；checksum 是结果摘要，用来发现输出变化或截断。
  \item[benchmark/profile/trace] benchmark 是受控性能测量；profile 是运行画像，说明时间或事件花在哪里；trace 是按时间保存的执行事件。三者都必须记录输入、构建选项和环境。
  \item[arena/workspace/batch/tensor] arena 是预先规划的一块内存区域；workspace 是临时工作区；batch 是成批处理的一组输入；tensor 是带形状、元素类型和步长的数据视图。
  \item[schema/metadata/manifest] schema 是结构规则；metadata 是描述数据的数据；manifest 是机器可读清单，记录构建、输入、输出、校验和和实验事实。
  \item[lane/mask/tail/stride] lane 是 SIMD 向量中的一个元素槽；mask 表示哪些 lane 有效；tail 是不能整除向量宽度或块大小的尾部；stride 是相邻元素在地址上的步幅。
  \item[handler/unwind/fallback] handler 是处理器或回调函数；unwind 是异常或调试器沿调用栈回退的过程；fallback 是主路径不可用时的保底实现。
  \item[AT\&T/GDT/FDE] AT\&T 是 GNU 工具链常见汇编语法风格，操作数顺序和 Intel 语法相反；GDT 是全局描述符表，保存段描述符；FDE 是 Frame Description Entry，DWARF 调用帧信息中的函数范围展开记录。
  \item[ITLB/DTLB/GC] ITLB 缓存取指地址翻译，DTLB 缓存数据地址翻译；GC 是垃圾回收。第一册提到 GC 时，多数是和 C/C++ 手动生命周期、碎片和移动式堆做对照。
  \item[ARM/GPU/GEMV] ARM 是一类处理器架构生态；GPU 是图形处理器或通用并行加速器；GEMV 是矩阵乘向量。它们在第一册里常作为后续平台或性能案例出现。
  \item[sanitizer/pass/pipeline] sanitizer 是编译器插桩诊断工具族；pass 是编译器或处理流水线中的一个分析/改写步骤；pipeline 是多个步骤按顺序组成的处理链。
  \item[hash/span/config/fixture] hash 是摘要或哈希值；span 是连续数据的非拥有视图；config 是配置；fixture 是测试用的固定输入、环境和预期上下文。
  \item[golden/contract/callback/worker/dispatcher] golden 是已确认的期望结果；contract 是边界约定；callback 是完成后被调用的函数；worker 是执行任务的线程或单元；dispatcher 是把任务分发到执行者的组件。
  \item[frontend/backend/dtype/shape/tokenizer] frontend 是把外部输入变成内部表示的前端；backend 是把内部表示落到机器或系统接口的后端；dtype 是元素类型；shape 是张量维度；tokenizer 是把文本切成 token 的组件。
  \item[microkernel/tile/packing/prefetch] microkernel 是极小的核心计算内核；tile 是分块计算单元；packing 是把数据重排成后端更容易连续读取的布局；prefetch 是提前把后续要用的数据拉近 CPU。
  \item[bandwidth/coherence/affinity/generation/owner] bandwidth 是单位时间数据搬运量；coherence 是多核缓存一致性；affinity 是线程或中断绑定到特定 CPU 的亲和性；generation 是版本编号；owner 是当前拥有对象或写权限的一方。
  \item[DIE/RSVD/IBT] DIE 是 DWARF Debugging Information Entry，调试信息里的条目；RSVD 是 reserved，保留位；IBT 是 Indirect Branch Tracking，CET 中限制间接跳转目标的一类控制流保护。
  \item[ModRM/GPR/RCX/RDX/R8/R9/XMM0] ModRM 是 x86 指令编码中的 mod-reg-r/m 字节；GPR 是 general-purpose register，通用寄存器；RCX、RDX、R8、R9 是通用寄存器名；XMM0 这类名字表示 128 位 SIMD 寄存器。寄存器名不是变量名，它们是 ABI 和指令实际读写的硬件状态。
  \item[MinGW/OpenMP/AoSoA] MinGW 是 Windows 上的 GNU 风格工具链；OpenMP 是用 pragma 表达并行循环和任务的编程模型；AoSoA 是 array-of-structs-of-arrays，介于结构数组和数组结构之间的数据布局。
  \item[CMake 变量] \code{DCMAKE_BUILD_TYPE} 选择 Debug/Release 等构建类型；\code{CMAKE_EXPORT_COMPILE_COMMANDS} 导出编译数据库；\code{CXXFLAGS} 和 \code{CFLAGS} 是 C++/C 编译选项。它们会改变生成代码和诊断证据。
  \item[Makefile 常量] \code{PHONY} 表示伪目标；\code{SRC}、\code{DEBUG_BIN}、\code{RELEASE_BIN} 常是项目自定义变量。它们不是系统概念，要回到当前构建脚本解释。
  \item[POSIX/ELF 宏] \code{RLIMIT_NOFILE}、\code{AT_FDCWD}、\code{WIFEXITED}、\code{WEXITSTATUS}、\code{WIFSIGNALED}、\code{WTERMSIG} 这类全大写名字是系统头文件或 ELF/进程接口里的宏。它们不是变量名，而是内核接口把资源限制、相对路径和子进程退出状态编码给用户态的方式。
  \item[benchmark 常量] \code{BENCHMARK_WARMUP_REPETITIONS}、\code{BENCHMARK_MEASURED_REPETITIONS}、\code{BENCHMARK_INPUT_SIZE}、\code{P95_NUMERATOR} 这类名字是实验代码的参数。它们负责把预热次数、测量次数、输入规模和分位数公式写清楚，避免报告只剩一串神秘数字。
  \item[汇编局部标签] \code{LJTI0_0}、\code{LBB0_1}、\code{FUNC}、\code{ENDP} 这类名字常来自汇编器、编译器或平台汇编语法。读它们时不要按 C++ 变量理解，要先判断它是跳转表、基本块标签、函数标记还是过程结束标记。
  \item[chunk/event/roofline] chunk 是连续切片或任务块；event 是一次可记录的状态变化或完成事实；roofline 是用算术强度、计算峰值和内存带宽估算性能上界的模型。
  \item[C++ 类型和转换] \code{size\_t} 表示对象大小和下标相关的无符号类型；\code{uintptr\_t} 是足以保存指针位模式的整数类型；\code{static\_cast} 是显式、受限的类型转换；\code{reinterpret\_cast} 是把同一批位按另一种类型解释的低层转换，必须非常谨慎。第一册看到这些名字时，要追问“这是数值转换、地址位模式，还是 ABI 边界”。
  \item[容量单位] KiB、MiB、GiB 明确按 1024 进位；KB、MB、GB 在工程材料里常按十进制或上下文约定使用。测量报告要写清口径，否则 cache、页大小、文件大小和带宽数字会被误读。
  \item[系统调用工具函数] \code{munmap} 解除 \code{mmap} 建立的映射；\code{clock\_gettime} 从内核或 vDSO 读取时钟；\code{sigaltstack} 设置信号备用栈。它们是用户态向系统边界发出的请求，不是普通 C++ 函数调用。
  \item[实验函数名] \code{sum\_u64}、\code{asm\_add2}、\code{bench\_sum}、\code{copy\_u64} 这类名字是本书实验函数：前缀 asm 常表示汇编实现，bench 常表示测量入口，u64/i32/f32 表示数据类型。读它们时先定位“输入类型、输出合同、测量目标”，再看机器码。
  \item[样本字段名] \code{elapsed\_ns}、\code{median\_ns}、\code{raw\_samples}、\code{bytes\_read} 这类字段是实验记录。ns 是纳秒，median 是中位数，raw samples 是原始样本；这些字段用于复查测量，不是硬件缩写。
  \item[vDSO] virtual Dynamic Shared Object，内核映射给用户态的一小段共享对象。某些系统调用风格的查询，例如读取时间，可以通过 vDSO 避免真正陷入内核，从而降低开销。
  \item[控制寄存器和中断名] CR3 是 x86 保存当前页表根地址的控制寄存器；IRQ 是 interrupt request，中断请求；IDT/GDT 分别是中断描述符表和全局描述符表。它们解释“CPU 如何找到内核入口和地址翻译上下文”，不是 C++ 变量。
  \item[链接器对象词] COMDAT 是目标文件中用于折叠重复定义的一类 section 机制；HIDDEN/GLOBAL 常表示符号可见性；ORIGIN 常见于动态库搜索路径规则。遇到这些词，要回到 ELF 符号、section 和动态链接器规则里解释。
  \item[教学短码] \code{B1/B2/B3}、\code{I1/I2/I3}、\code{S12/S24} 这类短码通常是图表或 trace 里的块、指令或样本编号。它们不是业界缩写，只是为了在一张图或一段实验里引用某个位置。
\end{description}
